\section{Preliminaries}

\textbf{Definition 1 (Urban Flow Observations).} For each city $c$ in a set $\mathcal{C}$ with $N_c$ spatial units, we observe a flow sequence $\mathbf{X}^{(c)}=[\mathbf{x}^{(c)}_1,\ldots,\mathbf{x}^{(c)}_{T_c}]\in\mathbb{R}^{T_c\times N_c}$. A forecasting sample $(\mathbf{X}^{(c)}_{t-H+1:t},\mathbf{Y}^{(c)}_{t+1:t+P})$ uses $H$ historical steps to predict the next $P$ steps, with $\mathbf{Y}^{(c)}_{t+1:t+P}=\mathbf{X}^{(c)}_{t+1:t+P}$.

\noindent\textbf{Definition 2 (Anomalous Flow Context).} Let $\mathbf{m}^{(c)}\in\{0,1\}^{T_c}$ be the anomaly context sequence, where $m^{(c)}_t=1$ marks a step affected by holidays, extreme weather, or other external disturbances. The future-window anomaly indicator is
\begin{equation}
    a^{(c)}_t=\mathbb{I}\left(\sum_{\tau=t+1}^{t+P}m^{(c)}_\tau>0\right).
\end{equation}
Samples with $a^{(c)}_t=1$ form the anomaly subset; all samples form the full subset.

\noindent\textbf{Problem Formulation (Cross-City Anomalous Urban Flow Forecasting).} Given a source city $c_s$ and a target city $c_t$ ($c_s\neq c_t$) that differ in node number, flow scale, periodicity, and anomaly distribution, we train on source samples $\mathcal{D}_s$ and directly apply the model to target test samples $\mathcal{D}_t$:
\begin{equation}
\begin{gathered}
    \theta^*=\arg\min_\theta\sum_{t\in\mathcal{T}_s}\ell\left(f_\theta(\mathbf{X}^{(c_s)}_{t-H+1:t}),\mathbf{Y}^{(c_s)}_{t+1:t+P}\right),\\
    \widehat{\mathbf{Y}}^{(c_t)}_{t+1:t+P}=f_{\theta^*}(\mathbf{X}^{(c_t)}_{t-H+1:t}).
\end{gathered}
\end{equation}
The key challenge is transferring forecasting knowledge to target-city anomalous periods under coexisting cross-city distribution shift and anomaly scarcity.